\documentclass[11pt]{amsart}
\usepackage[T1]{fontenc}
\usepackage[utf8]{inputenc}
\usepackage{amssymb}
\usepackage[margin=1.05in]{geometry}
\usepackage[hidelinks]{hyperref}
\newtheorem{theorem}{Theorem}[section]
\newtheorem{proposition}[theorem]{Proposition}
\newtheorem{lemma}[theorem]{Lemma}
\theoremstyle{definition}
\newtheorem{definition}[theorem]{Definition}
\theoremstyle{remark}
\newtheorem{remark}[theorem]{Remark}
\DeclareMathOperator{\Spec}{Spec}
\DeclareMathOperator{\im}{im}
\DeclareMathOperator{\Rees}{Rees}
\newcommand{\OO}{\mathcal O}
\newcommand{\FF}{\mathcal F}
\newcommand{\PP}{\mathbf P}
\newcommand{\code}[1]{\texttt{#1}}
\newcommand{\mere}[1]{\left\lVert #1\right\rVert}
\AddToHook{shipout/foreground}{%
 \put(0.5\paperwidth,-\paperheight+0.55cm){\makebox[0pt]{\tiny\ttfamily
 unit/doc/\jobname.tex\ ---\ compiled \today}}}
\title[Zariski Main Theorem by finite section lifting]{Zariski Main Theorem by finite section lifting\\
and the feasibility of a Unit formalisation}
\date{30 October 2026}
\begin{document}
\maketitle

\noindent\textbf{Scope and status.}
This note explains a projective proof of the version of Zariski Main
Theorem used on page~1 of \emph{Local structure of \'etale algebras}
(\code{zmt2-revised.pdf}). It first gives a classical proof using formal
functions, then isolates a finite constructive replacement. The latter
extends the method of \code{finite-lifting.tex} and the two-equation
ACL development. This document is a mathematical proof and a proposed
formalisation programme, not a new Unit-checked theorem. The assessment
below is based on the existing source interfaces and README files; no
new Unit declarations were written or checked for this note.

\medskip
\noindent\textbf{Method of simplification.}
The first task is to work backwards from the output certificate, not
to formalise the most general theorem appearing in the classical
proof. Section~\ref{sec:backwards} carries out this reduction. In
particular, it removes the need for a constructive theorem saying that
the entire ring of global functions is finite. The library programme
below should be read with this smaller target in mind. Further analysis
in Section~\ref{sec:sharper} reduces the remaining input to finite
top-degree coordinates from a bounded partial resolution. A relation
among those coordinates is enough; even the kernel of the associated
finite matrix need not be generated.

\section{The exact statement needed}

\begin{theorem}\label{thm:zmt}
Let $R$ be a commutative ring, $I\subseteq R$ an ideal, and
$E=R[x_1,\ldots,x_r]$ a finitely presented $R$-algebra. Suppose that
$E/IE$ is finite as an $R/I$-module. There exist a finite
$R$-subalgebra $B\subseteq E$ and $s\in B\cap(1+IE)$ such that the
inclusion induces an isomorphism
\[
 B[1/s]\xrightarrow{\ \sim\ }E[1/s].
\]
\end{theorem}

Here and throughout, ``finite'' means finitely generated as a module;
it does not mean free. A sufficient output certificate consists of
$s\in1+IE$ and monic equations over $R$ for
\[
 s,\quad sx_1,\ldots,sx_r.
\]
Indeed, $B=R[s,sx_1,\ldots,sx_r]$ is then finite, and
$x_i=(sx_i)/s$ proves the localisation assertion.

For the application to \code{zmt2-revised.pdf}, take $I=\mathfrak m$,
where $R$ is residually discrete local. The preceding finite-fibre
argument in that note supplies that $E/\mathfrak m E$ is finite.
Every prime $\mathfrak q$ containing $\mathfrak m E$ avoids
$s\in1+\mathfrak m E$. This supplies exactly the initial finite
subalgebra used there. No flatness, formal smoothness, or separability
is used in Theorem~\ref{thm:zmt} itself.

\section{Working backwards: a smaller certificate}\label{sec:backwards}

The finite algebra in Theorem~\ref{thm:zmt} is an output, but we need
not construct its multiplication table first. Nor do we need monic
equations for arbitrary global functions. The following smaller
certificate is sufficient, over any commutative base ring.

\begin{lemma}[One finite stable module]\label{lem:stable}
Let $V\subseteq E$ be a finitely generated $R$-submodule containing
$1$. Suppose multiplication by each of $s,sx_1,\ldots,sx_r$ preserves
$V$. Then all these elements are integral over $R$.
\end{lemma}
\begin{proof}
Choose generators $v_1,\ldots,v_m$ of $V$. For any one of the displayed
elements $t$, write $tv_j=\sum_k a_{kj}v_k$, with $a_{kj}\in R$.
The characteristic polynomial of $(a_{kj})$ is monic and annihilates
every $v_j$ by the adjugate identity. It annihilates $1\in V$, hence
vanishes at $t$ in $E$. Independence of the $v_j$ is not required.
The matrices can be different for the different $t$.
\end{proof}

Let $F_dE$ be the $R$-span of monomials of total degree at most $d$
in $x_1,\ldots,x_r$. It has an explicit finite generating list.

\begin{lemma}[Degree-lowering certificate]\label{lem:lower}
Suppose that $s\in1+IE$ and, for each $i$, there are $N_i\geq1$ and
$q_i\in R[T_1,\ldots,T_r]$ with
\[
 s x_i^{N_i}=q_i(x),\qquad \deg q_i\leq N_i-1.
 \tag{DL}
\]
Then Theorem~\ref{thm:zmt} holds for these data.
\end{lemma}
\begin{proof}
Choose a polynomial representing $s$, with degree bound $D\geq0$,
and put $d=D+\sum_i(N_i-1)$. Relation (DL) gives
\[
 sF_{d+1}E\subseteq F_dE.
\]
For a monomial $x^\alpha$ of degree at most $d+1$, if some
$\alpha_i\geq N_i$, replacement by (DL) gives
$s x^\alpha\in F_{|\alpha|-1}E\subseteq F_dE$. Otherwise
$|\alpha|\leq\sum_i(N_i-1)$, and the chosen polynomial for $s$
gives $s x^\alpha\in F_dE$ directly. Thus $F_dE$ contains $1$
and is stable under $s$ and all $sx_i$.
Lemma~\ref{lem:stable} proves the assertion.
\end{proof}

\begin{lemma}[Only chart denominators are needed]\label{lem:denom}
Let $W\subseteq\PP^r_R$ have affine chart $\Spec E=D_+(Z)$, and let
$h\in H^0(W,\OO_W(-1))$. Set $s=(Zh)|_{\Spec E}$. Then $s$ satisfies
the degree-lowering relations (DL).
\end{lemma}
\begin{proof}
On $D_+(X_i)$ represent $h$ as $Q_i/X_i^{k_i}$ with $k_i\geq1$
and $Q_i$ homogeneous of degree $k_i-1$ in the homogeneous
coordinate ring. Choose a homogeneous polynomial representative.
Agreement with $h=s/Z$ on $D_+(ZX_i)$ gives, after setting $Z=1$,
\[
 s x_i^{k_i}=Q_i(1,x)\quad\hbox{in }E[1/x_i].
\]
Clear this equality: for some $l_i\geq0$,
\[
 s x_i^{k_i+l_i}=x_i^{l_i}Q_i(1,x)\quad\hbox{in }E.
\]
Take $N_i=k_i+l_i$ and $q_i(T)=T_i^{l_i}Q_i(1,T)$, of degree
at most $N_i-1$. This uses only localisation equalities and degrees.
\end{proof}

Thus a lifted section gives the final certificate without computing
the whole ring $H^0(W,\OO_W)$, its generators, or its products. In
fact the extraction needs only agreement of $h$ on $D_+(Z)$ with
each $D_+(X_i)$; agreement between the other charts is not used.
Constructing a genuine global section remains one way to obtain this
weaker datum. We do not claim here that a lifting theorem for the
weaker overlap conditions has already been established.

\section{The projective proof over a noetherian ring}

For this section assume that $R$ is noetherian. Choose a presentation
\[
 E=R[T_1,\ldots,T_r]/(f_1,\ldots,f_t).
\]
Homogenise the $f_j$ in
$Z,X_1,\ldots,X_r$, and let $W\subseteq\PP^r_R$ be their common zero
scheme. Then $W$ is projective and
\[
 X:=W\cap D_+(Z)=\Spec E,\qquad x_i=X_i/Z.
\]
One need not compute the saturated ideal of the scheme-theoretic
closure: this projective scheme, possibly with extra components at
infinity, is sufficient. Put $\FF=\OO_W(-1)$.

\begin{lemma}\label{lem:nilpotent}
For every $n\geq1$, $E/I^nE$ is finite over $R/I^n$.
\end{lemma}
\begin{proof}
The image of each $x_i$ in the finite algebra $E/IE$ satisfies a monic
polynomial. Lift it to $P_i(T)\in R[T]$. Then $P_i(x_i)\in IE$, so
$P_i(x_i)^n\in I^nE$. The polynomial $P_i(T)^n$ is monic. Reducing
powers of each of the finitely many generators by these monic equations
gives a finite spanning family for $E/I^nE$. This argument also works
without noetherianity.
\end{proof}

Write $W_n=W\times_R R/I^n$ and $X_n=\Spec(E/I^nE)$.

\begin{lemma}\label{lem:sections}
The open subscheme $X_n\subseteq W_n$ is also closed. There is a
canonical section
\[
 h_n\in H^0(W,\FF/I^n\FF)
\]
equal to $Z^{-1}$ on $X_n$ and zero on its complement. These sections
are compatible as $n$ varies.
\end{lemma}
\begin{proof}
By Lemma~\ref{lem:nilpotent}, $X_n$ is finite, hence proper, over
$R/I^n$. Since $W_n$ is separated over that base, the map
$X_n\to W_n$ is proper (use its closed graph). Being also an open
immersion, it identifies $X_n$ with an open-and-closed subscheme.
The coordinate $Z$ trivialises $\OO(1)$ on $X_n$, so $Z^{-1}$ is a
section of $\FF$ there. Glue it with zero across the disjoint
complement. Both the decomposition and this section commute with
reduction from $n+1$ to $n$.
\end{proof}

This is the general replacement for the particular almost-section
$(0,0,Z^{-1})$ used for two residually linear equations. The finite
fibre need not consist of a single rational point.

\begin{proposition}\label{prop:formal}
There exists $h\in H^0(W,\FF)$ whose reduction modulo $I$ is $h_1$.
\end{proposition}
\begin{proof}
Set $M=H^0(W,\FF)$, a finite $R$-module by projective cohomological
finiteness. The theorem on formal functions gives
\[
 \widehat M\ \cong\ \varprojlim_n H^0(W,\FF/I^n\FF).
\]
Thus the $h_n$ determine $\widehat h\in\widehat M$. For a finite
module over a noetherian ring, $M/IM\cong\widehat M/I\widehat M$.
Choose $h\in M$ representing the reduction of $\widehat h$.
Naturality of the comparison maps gives the asserted reduction.
\end{proof}

\begin{remark}
We do not assert that $\widehat h$ comes from $M$. In general $R$ is
not complete, and that assertion would be false. Only its reduction
modulo $I$ must come from an actual section. Formal functions here is
the comparison theorem for a proper morphism, a coherent sheaf, and
an ideal of a noetherian base; see~\cite{Stacks}, Theorem~30.20.5.
\end{remark}

\begin{proof}[Conclusion of Theorem~\ref{thm:zmt} over noetherian $R$]
Set $s=(Zh)|_X$. Its reduction is $1$, because $h$ reduces to
$Z^{-1}$ on $X_1$. Lemma~\ref{lem:denom} gives degree-lowering
relations, and Lemma~\ref{lem:lower} produces monic equations for
$s,sx_1,\ldots,sx_r$. Take $B=R[s,sx_1,\ldots,sx_r]$.
It is finite and $B[1/s]=E[1/s]$. This proves exactly what is needed,
without a separate finiteness theorem for $H^0(W,\OO_W)$.
\end{proof}

\section{A finite replacement for formal functions}

Let $C^\bullet$ be the finite standard-chart \v Cech complex for $\FF$
on $W$. Its terms need not be finite $R$-modules. Sections of the
quotient on each affine intersection lift to sections of $\FF$;
consequently $h_n$ can be represented by $z_n\in C^0$ with
\[
 dz_n\in I^nC^1,\qquad z_n-z_1\in IC^0.
\]
The second assertion follows from compatibility modulo $I$.
Localisation is exact, so the subcomplex of sections of $I^n\FF$
is precisely $I^nC^\bullet$. None of these assertions assumes that
$\FF$ is flat over $R$.

\begin{proposition}[Finite lifting]\label{prop:finite}
Suppose the graded module
\[
 \mathcal H=\bigoplus_{n\geq0}H^1(I^nC^\bullet)T^n
\]
over $\mathcal R=\bigoplus_{n\geq0}I^nT^n$ has the following
constructive noetherian property: a sequence of homogeneous elements
has a term which is a finite linear combination of its predecessors,
with a certificate. Then there is $z\in C^0$ with
$dz=0$ and $z-z_1\in IC^0$.
\end{proposition}
\begin{proof}
For $n\geq1$ set $\beta_n=[dz_n]T^n\in\mathcal H_n$.
A dependence $\beta_N=\sum_{1\leq i<N}q_i\beta_i$, after taking
degree $N$, gives coefficients $r_i\in I^{N-i}$ and
$v\in I^NC^0$ with
\[
 dz_N-\sum_{1\leq i<N}r_i dz_i=dv.
\]
Then
\[
 z=z_N-\sum_{1\leq i<N}r_i z_i-v
\]
is closed. Since every $r_i$ is in $I$, it has the desired reduction.
The case $N=1$ has an empty sum. Only finitely many approximations
and one finite relation occur in the certificate.
\end{proof}

Classically, $\mathcal H$ is finite over the noetherian Rees algebra
by projective cohomological finiteness applied over that algebra;
this is also Lemma~30.20.1 of~\cite{Stacks}. Constructively, one must
produce the generators and the noetherian dependence mechanism.
These are explicit obligations, rather than an appeal to completion.
This proposition is the argument of \code{finite-lifting.tex}.

Under propositional truncation, a family of merely existing lifts
must not be replaced by a chosen infinite sequence. Instead, given
a coded noetherian bar, traverse its well-founded tree: at each node
use one truncated lift and code, recurse into the proposition-valued
goal, and stop at a certified dependence. The existing
\code{GradedDependence} development is a model for this discipline;
the generic construction above is not already its exported interface.

\subsection*{Relation to Artin--Rees and to flatness}
Artin--Rees controls the induced filtration on a submodule:
\[
 N\cap I^{m+c}R^a\subseteq I^mN.
\]
Our existing two-equation proof uses this to express a small
removable obstruction with small coefficients on fixed removable
obstructions. Their primitives then give a small correction of a
cochain. Proposition~\ref{prop:finite} packages the analogous control
in a graded module over the Rees algebra. One can pursue either
implementation; using the graded proposition does not require first
formalising the entire theorem on formal functions.

The flatness argument uses a different consequence of Artin--Rees,
$\operatorname{Ann}(f)\cap f^NE=0$ for sufficiently large $N$.
Finite Taylor lifts place an annihilated element deep enough in this
filtration to kill it. This explains the common algebraic mechanism,
but flatness is not an input to the present ZMT proof. In particular,
the Rees sheaf must use the images $I^n\FF$, not replace them by
$I^n\otimes_R\FF$ without justification.

\section{Reduction to an arbitrary base ring}\label{sec:descent}

Return to arbitrary $R$. From finiteness of $E/IE$, choose monic
$P_i\in R[T]$ and finite witnesses
\[
 P_i(x_i)=\sum_{j=1}^u a_j H_{ij}(x),\qquad a_j\in I.
\]
A single finite family $a_j$ suffices for all $i$. Choose polynomial
witnesses for these identities in the presentation of $E$; thus each
difference is expressed as a combination of the $f_l$.

Introduce finitely many variables over $\mathbb Z$ for all coefficients
in the presentation, the $P_i$, the $H_{ij}$, the $a_j$, and the
identity witnesses. Impose the finitely many coefficient equalities
which make the displayed identities hold. This constructs a finitely
presented $\mathbb Z$-algebra $R_0$, an ideal
$I_0=(a_{1,0},\ldots,a_{u,0})$, and a finitely presented algebra $E_0$
with a map to the original data. Monicity is preserved by choosing
the leading coefficients to be $1$. The equations imply that
$E_0/I_0E_0$ is finite, by reduction of powers.

Apply the noetherian-base result to $(R_0,I_0,E_0)$. Transport its
finite certificate to $E$. More explicitly, if $B_0\subseteq E_0$
is finite, take the image of $R\otimes_{R_0}B_0$ in $E$.
This is finite, contains $s,sx_i$, and has $s\in1+IE$.
The map $R_0\to R$ need not be injective. One must not instead assume
that its image is a coherent ring or has a finitely generated kernel.
This reduction does not require $I$ itself to be finitely generated.

\section{What the existing Unit developments provide}

The following inventory records interfaces inspected for this note,
not a fresh validation run of those developments. Paths are relative
to \code{unit/sag/}.

\begin{itemize}
\item \code{artin-rees/CFN}, \code{ZCFN}, \code{HilbertCFN},
\code{QuotCFN}: coherent, strongly discrete, constructively noetherian
rings, with the integer, polynomial and finite-quotient instances.
These are suitable for the $R_0$ of Section~\ref{sec:descent}.
\item \code{artin-rees/FreeCoherent}, \code{FreeDisc},
\code{Saturation}, \code{DegreeCor}, \code{Schreyer}:
finite syzygies, bounded-degree intersections and polynomial
coherence. These are ingredients for general elimination and graded
presentations; they are not already a projective cohomology API.
\item \code{acl-zmt/Laurent3}, \code{Cech3},
\code{CechContraction}, \code{CechMonomials},
\code{CechH1Finite}: explicit three-coordinate Laurent and \v Cech
calculations. Their interfaces concern $\PP^2$, not arbitrary $\PP^r$.
\item \code{acl-zmt/FiniteCorrection}, \code{GradedDependence},
\code{ReesLiftingCong}: the finite correction and graded dependence
mechanisms, in interfaces using two ideal generators and particular
finite-coordinate obstruction data.
\item \code{artin-rees/ArtinRees}, \code{PowSpan},
\code{ACLRelations}, \code{ACLCorrection}: the alternative correction
by Artin--Rees for $I=(a,b)$, with the two-equation obstruction matrix.
The general number of ideal generators is still an extension.
\item \code{acl-zmt/IntegralSections} and the exterior-algebra
determinant files provide models for the passage from finite sections
to monic equations. \code{ACLCertificateGeneral} assembles the current
two-equation case, with its stated gcd-domain and CFN hypotheses.
\end{itemize}

The directory \code{sag/projective} constructs particular projective
line examples and coordinate calculations, with additional SAG
hypotheses in some results. Its name does not indicate a library of
general projective cohomological finiteness. The proposed proof can
stay with algebraic localisations and finite \v Cech complexes and
need not import those extra hypotheses.

\section{Further reductions of the remaining lifting problem}\label{sec:sharper}

\subsection{The one-variable case has a direct certificate}
Suppose $E=R[x]$. Lift a residual monic relation of degree $d\geq1$
and its $I$-membership witnesses to a polynomial relation
\[
 f(x)=0,\qquad f(T)\equiv P(T)\pmod{I R[T]},\qquad P\text{ monic},
 \quad\deg P=d.
\]
Writing $f(T)=\sum_{j=0}^m c_jT^j$ with $m\geq d$ (padding with
zero coefficients if necessary), put
\[
 s=c_d+c_{d+1}x+\cdots+c_mx^{m-d}.
\]
Then $s\in1+IE$ and
\[
 sx^d=-\sum_{j<d}c_jx^j.
\]
This is exactly (DL); Lemma~\ref{lem:lower} finishes without any
noetherianity, formal lifting, Artin--Rees, or cohomology. If a
residual relation has degree zero, enlarge it to positive degree.
This is a useful base case and an explicit test of what the
projective machinery must accomplish in more variables.

\subsection{The apparent weaker-overlap shortcut has a cost}
Section~\ref{sec:backwards} observes that extraction only uses
overlaps with $D_+(Z)$. It does not follow that one should simply
discard the other overlaps from the lifting complex. For example,
on $\PP^2_k$, retain just the edges $XZ$ and $YZ$. In degree $-1$,
the edge cochains
\[
 (X^{-n}Y^{2n-1}Z^{-n},\ 0),\qquad n\geq1,
\]
give linearly independent classes modulo chart coboundaries. A chart
$X$ cochain has nonnegative $Z$ exponent; a chart $Z$ cochain has
nonnegative $X$ exponent; neither can contribute the displayed
monomial on $XZ$. The chart $Y$ does not contribute on that edge.
Thus this shortened complex has infinite-dimensional cokernel,
although the actual $H^1(\PP^2_k,\OO(-1))$ is zero. Keeping only
the conditions used at the end can destroy the finiteness that makes
lifting possible. A different argument for the weaker conditions
would have to address this explicitly.

\subsection{Replace general cohomology finiteness by one finite matrix}
Here is a more effective reduction. Let $D$ be CFN,
$S=D[X_0,\ldots,X_r]$, and $M=S/J$, with $J$ a finitely generated
homogeneous ideal. Write all the following modules with their twist
by $-1$ already included. Choose a partial homogeneous free resolution
\[
 F_{r-1}\xrightarrow{\partial_{r-1}}F_{r-2}\longrightarrow
 \cdots\longrightarrow F_1\longrightarrow F_0=S(-1)
 \longrightarrow M(-1)\longrightarrow0.
 \tag{PR}
\]
It must be exact at $F_0,\ldots,F_{r-2}$; no condition on the kernel
at the left end is imposed. Each $F_j$ is a finite sum of shifts of
$S$. For $r=2$, $F_1$ is just a generating list for $J$, shifted.
For larger $r$, only $r-2$ additional rounds of homogeneous syzygies
are required after that generating list.

\begin{proposition}[A finite source for the needed classes]\label{prop:matrix}
For $r\geq2$, form the finite matrix
\[
 A:H^r(\widetilde F_{r-1})\longrightarrow
       H^r(\widetilde F_{r-2})
\]
induced by $\partial_{r-1}$. There is a surjection
\[
 \ker A\ \twoheadrightarrow\ H^1(\widetilde M(-1)).
\]
The source of $A$ and its target are finite free $D$-modules, with
the usual all-negative-exponent monomial bases for each shift.
Consequently coherence of $D$ supplies a finite generating family
for the required $H^1$ classes.
\end{proposition}
\begin{proof}
Put $K_0=M(-1)$ and use the exact sequences
$0\to K_{j+1}\to F_j\to K_j\to0$ through $j=r-2$,
with $F_{r-1}\twoheadrightarrow K_{r-1}$.
Intermediate cohomology of every shifted free module vanishes.
Also $H^q(\OO(-1))=0$ for all $q$, so the first connecting map
gives $H^1(K_0)\cong H^2(K_1)$. Continuing the connecting maps gives,
at the last step, an identification of $H^1(K_0)$ with
\[
 \ker\bigl(H^r(K_{r-1})\to H^r(F_{r-2})\bigr).
\]
For $r=2$ this is the first step itself. Since there is no cohomology
above degree $r$, the map
$H^r(F_{r-1})\to H^r(K_{r-1})$ is surjective. The preimage of the
displayed kernel is precisely $\ker A$. This proves the assertion.

The maps can be realised as finite \v Cech zigzags: choose Laurent
representatives in top degree, use $A v=0$ to remove the top
obstruction by the free-module monomial contraction, and descend
through (PR). Conversely, lift a given $1$-cocycle through (PR) and
apply these contractions to obtain a vector in $\ker A$. A formal
proof must implement both the zigzag and its equality witnesses;
the long exact sequence argument identifies exactly what to implement.
\end{proof}

For $r=2$, the target is $H^2(\OO(-1))=0$: every top-degree vector
is admissible. This recovers the special shape of
\code{CechH1Finite}. For $r=3$, one syzygy module and one finite
kernel calculation suffice. For $r=1$, $H^1(\widetilde M(-1))=0$
directly, since it is a quotient of $H^1(\OO(-1))$ and the complex
has no degree $2$. The direct polynomial argument above is preferable
for this base case.

\begin{remark}[A concrete first test in four coordinates]
Take $J=(X_0^3,X_1^3)$ in $D[X_0,X_1,X_2,X_3]$. The pair of
variable powers is regular, so its two-step Koszul resolution gives
\[
 F_2=\OO(-7),\qquad F_1=\OO(-4)^2,\qquad F_0=\OO(-1),
 \qquad\partial_2=(-X_1^3,X_0^3).
\]
Here the notation denotes the associated sheaves of the shifted
free modules. The source of $A$ has the $20$ monomials
$X_0^{-a_0}\cdots X_3^{-a_3}$ with $a_i\geq1$ and
$\sum a_i=7$; the target has one all-negative monomial in each
of its two summands. Its first row selects, with a minus sign,
$(a_0,a_1,a_2,a_3)=(1,4,1,1)$; its second selects $(4,1,1,1)$.
All other columns vanish. Thus its kernel is freely generated by
the remaining $18$ monomials, over any $D$. The zigzag must turn
these into $18$ generating $1$-cocycles. This tests a genuinely
nonzero top matrix without first requiring an arbitrary syzygy
algorithm. It is a test of the cohomology reduction, not an example
of a finite affine fibre. No Unit check of this example is claimed.
\end{remark}

\subsection{A further shortcut: do not generate the finite kernel}
The last sentence of Proposition~\ref{prop:matrix} still proves more
than we need. Let $V=H^r(\widetilde F_{r-1})$, a fixed finite free
$D$-module. For each particular obstruction class $\beta_n$, the
surjectivity argument gives a vector $v_n\in V$ with $Av_n=0$ and
image $\beta_n$. Any $D$-linear relation among these chosen vectors
is a relation among their images. Noetherian dependence can be
applied in $V$ itself. There is no need to find generators of
$\ker A$, or even to construct $A$ as a separate matrix interface.
Its kernel condition describes the domain of the connecting map;
the given zigzag certificates already supply that condition.

The useful formal interface is consequently smaller than
Proposition~\ref{prop:matrix}: send a given $1$-cocycle to finite
top-degree coordinates with a zigzag certificate, and prove that a
relation among finitely many such coordinates produces a primitive
for the same relation among their cocycles. The coordinate assignment
need not be canonical or additive. What must be additive is the
connecting construction on the certified top cycles. Constructive
noetherianity is used only in the ambient finite free module.

\begin{lemma}[The minimal relation-certificate interface]\label{lem:interface}
Let $D$ be CFN, $Z$ a $D$-module, $B\subseteq Z$ a submodule, and
$V=D^v$. Suppose a relation $\Gamma\subseteq Z\times V$ contains
$(0,0)$ and is closed under addition, negation and $D$-scalars.
Suppose also
\[
 \Gamma(z,0)\ \Longrightarrow\ z\in B.
\]
Let $b_1,b_2,\ldots$ be given elements of $Z$, each merely admitting
$p_n\in V$ with $\Gamma(b_n,p_n)$. Then there merely exist $N\geq1$
and $Q_i\in D$ for $1\leq i<N$ with
\[
 b_N-\sum_{i<N}Q_i b_i\in B.
\]
\end{lemma}
\begin{proof}
Use noetherian free-list induction in $V$, retaining in a finite
prefix the chosen $p_i$ and their certificates. At the next index,
eliminate the truncated existence of one further $p_n$ into the
proposition-valued goal. If it is in the span of its predecessors,
extract a finite dependence; closure of $\Gamma$ gives
$\Gamma(b_N-\sum Q_i b_i,0)$, hence the conclusion. Otherwise extend
the free list and recurse. This uses the same induction as
\code{SeqDependence}, with the choices inside the induction rather
than a chosen infinite sequence. The weighted homogeneous version
retains the degree witnesses and takes the appropriate homogeneous
parts of the coefficients in the resulting finite relation.
\end{proof}

In the application, $Z$ consists of the relevant graded cocycles,
$B$ consists of coboundaries, and $\Gamma$ is the existence of a
finite zigzag certificate. The lemma only needs module predicates
and finite lists, not quotients of $Z$ or a category of complexes.
It makes the remaining implementation obligation precise: construct
the particular zigzag relation and prove its closure and zero-vector
property. No further mathematical shortcut past this obligation is
established in this note.

\subsection{Use polynomial coordinates instead of building a Rees API}
Apply the preceding reduction to the bigraded presentation of
$N=\bigoplus_n I_0^n(S/J)T^n$ described below, over
$D=R_0[U_1,\ldots,U_u]$. Keep projective degree and Rees degree
separate. The top monomial lists are finite over $D$ and have fixed
Rees shifts $e_1,\ldots,e_v$. The raising zigzag for $[dz_n]T^n$
gives top-degree vectors $p_n=(p_{n,j}(U))_j$, with
$\deg_U p_{n,j}=n-e_j$ (zero when this is negative), together with
zigzag witnesses. We only need a dependence
among these vectors in the finite free $D$-module $D^v$.
Constructive noetherian dependence there gives
\[
 p_N=\sum_{1\leq i<N}Q_i(U)p_i.
\]
Taking weighted degree $N$, with the $j$th coordinate shifted by
$e_j$, replaces $Q_i$ by its homogeneous part of degree $N-i$.
Substitute $U_j=a_j$: its value $r_i$ belongs to $I_0^{N-i}$.
The relation-reflection property of the zigzag gives precisely
\[
 dz_N-\sum_{i<N}r_i dz_i=dv,\qquad v\in I_0^NC^0,
\]
and finite correction finishes. Terms of negative coordinate degree
are treated as zero, rather than assigned fictitious polynomials.

Thus no general cohomology quotient construction or standalone Rees
ring library is needed for the dependence step. One needs a
bigraded polynomial presentation, a bounded partial resolution,
finite top-coordinate lists, and the certificate that their relations
give coboundaries. Generators for the cohomology or the top matrix
kernel are unnecessary.
The noetherian-tree treatment of truncated choices remains necessary.
This is closely matched to the existing finite-coordinate
\code{GradedDependence} mechanism.

\section{The new work and its feasibility}

\subsection{Finite-fibre sections in an algebraic chart model}
Lemma~\ref{lem:nilpotent} and the final certificate are short algebraic
tasks. For Lemma~\ref{lem:sections}, one can avoid formalising the
general properness argument. Here is an algebraic construction of its
chart idempotents.

Work modulo $I^n$ on the chart $D_+(X_i)$, with coordinate ring $A_i$,
and put $u=Z/X_i$. On its intersection with $D_+(Z)$ we have
$x_i=u^{-1}$. A monic relation of degree $d$ for $x_i$ gives
\[
 H(u)=1+c_1u+\cdots+c_du^d=0\quad\hbox{in }A_i[1/u].
\]
Equality in a localisation supplies $L\geq1$ with $u^LH(u)=0$ in
$A_i$. Since $H(0)=1$, there are polynomials $A(U),B(U)$ with
\[
 A(u)u^L+B(u)H(u)=1.
\]
They can be obtained by the finite geometric-series inverse of $H$
modulo $U^L$. The element $e=A(u)u^L$ is idempotent, and
$D(e)=D(u)$. On the $e$-part, $u$ has inverse
$t=eA(u)u^{L-1}$, extended by zero on the complement; indeed
$ut=e$ and $et=t$. The local section is $tX_i^{-1}$, whereas on
$D_+(Z)$ it is $Z^{-1}$. These agree on overlaps because both are
$Z^{-1}$ on the affine part and zero on its complementary idempotent
part. The same characterisation proves compatibility under reduction.

Thus the finite-fibre section construction has explicit polynomial
certificates. The new formal work is to implement these formulas,
their chart changes and their overlap identities. The existing
\code{FixedCochain} treats only the residual-point special case.

\subsection{Any finite number of projective coordinates}
Define Laurent polynomials on $\mathbb Z^{r+1}$ with finite support,
partial localisation support conditions, and the alternating complex
indexed by nonempty subsets of the coordinate set. Generalise the
monomial contraction: the negative coordinates of a monomial determine
its possible chart components. This gives vanishing in intermediate
degrees for $\OO(k)$ and a finite monomial basis in top degree, as
well as the degree-zero calculation. Coefficients need not have
decidable equality for this contraction; the exponent comparisons are
decidable. This is a substantial generalisation of the three-variable
code, with a clear combinatorial proof.

\subsection{Finite presentations of the cohomology}
For our end result, the finite top-coordinate certificates of
Section~\ref{sec:sharper} suffice. A general cohomology-presentation
interface would do extra work. No $H^0$ calculation is needed for
the extraction of monic equations: Section~\ref{sec:backwards}
already supplies a smaller argument.

For a finitely presented graded module over
$D[X_0,\ldots,X_r]$, with $D$ CFN, construct successive homogeneous
syzygies and finite sums of shifted free modules mapping onto them.
Proposition~\ref{prop:matrix} explains why terms through $F_{r-1}$
suffice. Its further reduction avoids computing the finite kernel:
lift only the given obstruction classes to top-coordinate vectors,
and use dependence in their ambient finite free module. Coherence
is still used to construct the successive homogeneous syzygies over
the polynomial ring, not to present the resulting cohomology.

This is an implementation strategy, not an existing Unit theorem.
The truncation bound, exactness, lifts and comparison maps must be
proved. In particular, one must not assume that an arbitrary module
over a singular base has a terminating finite free resolution.
No general Hilbert syzygy theorem is required for a bounded initial
resolution. The convenient two-equation gcd syzygy calculation is
replaced here by iterated homogeneous coherence.

\subsection{Presenting the Rees module without a flatness assumption}
Let $S=R_0[X_0,\ldots,X_r]$, $J$ the chosen homogeneous ideal,
$M=S/J$, and $I_0=(a_1,\ldots,a_u)$. The relevant graded algebra is
\[
 N=\bigoplus_{n\geq0} I_0^nM\,T^n\subseteq M[T].
\]
Set $D=R_0[U_1,\ldots,U_u]$. The map
\[
 D[X_0,\ldots,X_r]\longrightarrow M[T],\qquad U_j\longmapsto a_jT
\]
has image $N$. Its kernel is the elimination ideal
\[
 \bigl(J,\ U_1-a_1T,\ldots,U_u-a_uT\bigr)
 R_0[U,X,T]\ \cap\ R_0[U,X].
\]
It is not in general enough to impose $J$ and the familiar pairwise
Rees relations: $M$ was not assumed flat over $R_0$.

There is a useful potential reuse here. Treat $T$ as the final
polynomial variable over the CFN ring $R_0[U,X]$. The degree-zero
intersection case of \code{DegreeCor} supplies generators of precisely
this elimination ideal, with membership witnesses. Thus a new
multivariate Gr\"obner-basis implementation need not be a prerequisite.
One still has to implement the bridge to that interface, extract
homogeneous components for the projective and Rees gradings, and prove
the kernel/image presentation. This proposal has not been instantiated
in Unit for the present ideal.

Applying the finite-coordinate construction over $D$ to this
presentation controls the required classes in
\[
 \bigoplus_n H^1(W,I_0^n\OO_W(-1)).
\]
One must prove that localisation, the finite chart complex and
extraction of projective degree commute with the direct sum. Keep
polynomial coordinates and use dependence in $D^v$, as in
Section~\ref{sec:sharper}. There is no need to construct a separate
formal Rees-ring interface or generate the entire graded cohomology
module. The two gradings and their comparison maps remain a
significant source of representation work.

\subsection{Assessment and suggested milestones}
The proof is feasible in the foundations and algebraic style already
used by Unit: finite supports, quotients, images, propositional
truncations and coded noetherian trees. There is no identified need
for countable choice, excluded middle, flatness, or a full
formal-functions library. Nevertheless it is a substantial library
extension, not a short instantiation of the existing ACL theorem.
The sharper target is the free-module monomial contraction and the
finite-coordinate zigzag explained by Proposition~\ref{prop:matrix}, together with
their bigraded presentation bridges. A general projective cohomology
library is optional. The polynomial-coordinate route also avoids
generalising Artin--Rees itself to arbitrarily many ideal generators.
Only finite ideal powers and elementary correction formulas need
that generalisation. The finite-fibre section construction remains
a separate task with explicit chart formulas.

A useful order of work, revised by working backwards, is:
\begin{enumerate}
\item Formalise Lemmas~\ref{lem:stable}--\ref{lem:denom}: extract
the certificate using bounded-degree monomials and chart
denominators. Then formalise the finite-presentation descent of
Section~\ref{sec:descent}. Neither task needs cohomological finiteness.
\item Generalise ideal powers and finite correction to a finite list
of generators. Use noetherian dependence over $D$ to produce the
weighted polynomial-coordinate relation. State the lifting result
conditionally on its explicit finite-coordinate reduction data.
\item Implement finite-fibre sections and the Laurent/\v Cech
contraction for any fixed finite number of coordinates. Check the
new construction against the existing $\PP^2$ formulas.
\item Implement the raising \v Cech zigzag and its relation-reflection
certificate, initially for a supplied partial homogeneous resolution.
Then construct the terms through $F_{r-1}$ using CFN coherence.
Do not compute a top matrix kernel or cohomology generators unless
another application needs them.
\item Instantiate elimination for $N$, verify the Rees/\v Cech
comparison, and discharge the graded finiteness hypothesis. Assemble
the arbitrary-base theorem and the instance $I=\mathfrak m$.
\end{enumerate}

Each milestone should export a theorem with its genuine parameters.
A conditional lifting result is useful, but does not prove ZMT until
the finiteness and finite-fibre inputs are constructed. The first
decisive feasibility check beyond the existing $\PP^2$ code is a
$\PP^3$ example: one homogeneous syzygy module, the $20$ top
monomials in the cubic example above, a raising zigzag, and the
primitive obtained from a relation among its top-coordinate vectors.
No reliable calendar estimate follows from the present inspection.

\subsection{Current stopping point and practical feasibility}
The simplifications above leave a concrete finite-algebra programme.
We do not have a proved route which dispenses with the uniform
finite-coordinate construction. Merely imposing fewer overlap
conditions fails by the infinite-dimensional example above; merely
applying noetherianity to the original chart modules is unavailable,
since those modules are not finite over the base. These observations
identify the remaining difficulty, not an impossibility theorem for
all other approaches.

The existing determinant trick, polynomial coherence, saturation,
finite free-module dependence and finite correction provide the
algebraic ends of the proof. The principal new proof is the zigzag
relation of Lemma~\ref{lem:interface}, for general finite numbers of
coordinates and a supplied partial resolution. Weighted homogeneous
syzygies, the elimination presentation and the comparison with
$I^n$-images then connect it to the application. Existing coherence
is ungraded: extracting and tracking the two homogeneous gradings is
new interface work, not an already verified automatic operation.

A prudent first implementation should therefore stop at two precise
deliverables: the abstract relation-certificate lemma, including
truncated choices; and the four-coordinate cubic example with its
zigzag and relation-to-primitive theorem. Completing those would
validate the new mechanism. Completing only the abstract lemma would
not validate projective lifting. The full theorem is now a realistic
but substantial formalisation, with representation and comparison
proofs as the main risks. There is no basis yet for calling it a
small patch or assigning a reliable completion time. A machine-checked
mathematical certificate is the target; an efficient executable ZMT
algorithm would be a further requirement.

\section{Checks which should accompany a formalisation}

The output should include $s-1=\sum_j a_jr_j$ and explicit monic
equations for $s,sx_i$, with polynomial ideal-membership witnesses in
the presentation of $E$. Finite generation rather than freeness must
be used throughout. Useful regression cases are:
\begin{itemize}
\item the existing two-variable ACL example, to recover the current
certificate theorem;
\item $E=R[t]/(t^2)$ and $I=0$, to ensure that a nonreduced finite
fibre is allowed;
\item $R=k[a]$, $I=(a)$, $E=R\times R[1/a]$, with the diagonal
$R$-algebra structure: $B=R\times R$ and $s=(1,a)\in1+IE$ work.
The affine fibre is one component; a component invisible modulo $I$
must not be lost before localisation;
\item $E=R[1/a]$, $I=(a)$, whose special fibre is zero. A valid
certificate may have $s=0$, since $IE=E$, giving zero localisations;
\item a nonflat finite algebra over $R$, to detect accidental use of
$I^n\otimes\FF=I^n\FF$.
\end{itemize}

The proposed new Unit files should follow the proof one lemma per
declaration, with explicit structures and named comparison maps as
in \code{seg/WRITING.md}. Validation must include fresh Unit checks
of those files and their exact exported assumptions. Compiling this
LaTeX note only validates the document, not the proposed formalisation.

\begin{thebibliography}{9}
\bibitem{Horrocks} G.~Horrocks, \emph{Projective modules over an
extension of a local ring}, Proc. London Math. Soc. (3) 14 (1964),
714--718; local file \code{phorrocks.pdf}. The section lifting in
Theorem~1 is the motivation for this note.
\bibitem{Local} T.~Coquand, \emph{Local structure of \'etale algebras},
local file \code{/Users/coquand/COLL/zmt2-revised.pdf}, especially p.~1.
\bibitem{ACL} M.~Alonso, T.~Coquand and H.~Lombardi,
\emph{Revisiting Zariski Main Theorem from a constructive point of
view}, Journal of Algebra 406 (2014), 46--68.
\bibitem{Stacks} The Stacks Project,
\href{https://stacks.math.columbia.edu/tag/02O7}{Section 30.20,
The theorem on formal functions}, especially Lemma~30.20.1 and
Theorem~30.20.5. The proofs above specify the particular consequence
needed rather than claiming a new proof of the entire theorem.
\bibitem{Internal} Internal notes in \code{unit/doc/}: the files
\code{finite-lifting.tex} and \code{finite-cohomology.tex}; also
\code{artin-rees.tex} and \code{artin-rees-acl-elementary.tex}.
Their historical status paragraphs should be read with the current
source and README files.
\end{thebibliography}
\end{document}
